% =============================================================
\section{Discussion and Limitations of the Compton Formula}
% =============================================================

The derivation and experimental confirmation of the Compton formula
\[
\Delta \lambda = \frac{h}{m_e c}(1 - \cos\theta)
\]
was a milestone in the history of physics. Nevertheless, it is valid only within certain limits
and is based on idealized assumptions. This chapter discusses its range of validity,
its limitations, and its significance within the context of modern theories.

% -------------------------------------------------------------
\subsection{Range of Validity}
% -------------------------------------------------------------

The Compton formula describes the elastic scattering of a photon by an electron under the assumptions that
\begin{itemize}
	\item the electron is initially at rest,
	\item the collision is elastic (\(E_\text{photon} + E_\text{electron} = \text{constant}\)),
	\item and only \emph{one photon} interacts with \emph{one electron}.
\end{itemize}

\begin{PhysicsBox}[Validity of the Compton Formula]
	The formula is strictly valid only for free electrons.
	In real materials, however, electrons are bound to atoms.
	Their binding energies cause small deviations from the ideal curve,
	which are observed experimentally as the \emph{Compton profile}.
\end{PhysicsBox}

% -------------------------------------------------------------
\subsection{Limitations at High Photon Energies}
% -------------------------------------------------------------

For very short-wavelength photons, especially in the gamma range (\(E_\gamma \gtrsim 1~\text{MeV}\)),
the Compton effect is accompanied by additional quantum-mechanical processes:

\begin{itemize}
	\item The recoil of the electron becomes relativistic.
	\item For energies above \(1.022~\text{MeV}\), pair production (\(\gamma \to e^+ + e^-\)) can occur.
	\item The scattering cross section becomes strongly angle-dependent and is described by the \emph{Klein–Nishina formula}.
\end{itemize}

\begin{DidacticBox}[Relativistic Extension]
	For very high-energy photons, the classical Compton formula
	must be replaced by the \emph{Klein–Nishina equation}.
	This equation describes the differential scattering cross section
	in a relativistically correct way and originates from quantum electrodynamics (QED).
\end{DidacticBox}

% -------------------------------------------------------------
\subsection{Limitations at Low Photon Energies}
% -------------------------------------------------------------

For long-wavelength photons (\(h\nu \ll m_e c^2\)), the Compton shift becomes so small
that it can no longer be detected experimentally. In this limit,
the Compton effect reduces to classical \emph{Thomson scattering}.

\begin{NoteBox}[Transition to Thomson Scattering]
	In the low-energy regime, photons behave like classical electromagnetic waves.
	The electrons oscillate in the light field and re-radiate—
	the wavelength remains practically unchanged.
\end{NoteBox}

% -------------------------------------------------------------
\subsection{Theoretical Significance and Model Limitations}
% -------------------------------------------------------------

The Compton formula provides compelling evidence for the particle nature of light,
but it still treats the photon and the electron as classical collision partners.
In modern quantum field theory, scattering is understood as an interaction
between two quantized fields — the electromagnetic field
and the electron field.

\begin{DidacticBox}[Limit of the Wave–Particle Picture]
	The Compton effect can be intuitively explained using photon momentum,
	but the concept of the “photon as a particle” is only an approximation.
	In quantum electrodynamics, the effect arises from the exchange of virtual photons
	between the electron and the electromagnetic field.
\end{DidacticBox}

% -------------------------------------------------------------
\subsection{Historical Significance}
% -------------------------------------------------------------

\begin{HistoryBox}[Compton’s Legacy]
	Arthur H. Compton received the 1927 Nobel Prize in Physics
	for demonstrating experimentally that light carries quantum momentum.
	His work confirmed the dual nature of light and marked
	the transition from classical wave theory to quantum theory.
\end{HistoryBox}

% -------------------------------------------------------------
\subsection*{Summary}
% -------------------------------------------------------------

\begin{itemize}
	\item The Compton formula is exact for free electrons undergoing elastic scattering.
	\item For very high energies, relativistic corrections (Klein–Nishina) must be taken into account.
	\item For low energies, the shift vanishes — leaving only Thomson scattering.
	\item The formula marks the transition from classical optics to quantum physics.
\end{itemize}

\begin{PhysicsBox}[Significance for Modern Physics]
	The Compton effect unites light, matter, and momentum conservation in a single relationship.
	It was the first direct experimental proof of photon momentum
	and still forms a bridge between classical electrodynamics
	and quantum electrodynamics.
\end{PhysicsBox}
